\chapter{Summary and conclusions}\label{chap:summary_conclusions}
Coherent Elastic Neutrino-Nucleus Scattering (CE$\nu$NS) provides a uniquely clean laboratory for probing the electroweak sector of the Standard Model (SM) at low momentum transfer, since its cross section factorizes into a coherently enhanced nuclear form factor and a weak charge that carries the full dependence on the underlying couplings. In this work we exploited that sensitivity to study two related, and often degenerate, SM quantities: the weak mixing angle $\sin^2\theta_W$ and the electron neutrino charge radius (NCR) $\langle r_{\nu_e}^2\rangle$, using a liquid argon (LAr) detector exposed to the antineutrino flux of a research reactor. We began by tracing the historical and formal path that leads from Fermi's effective four-fermion interaction, through the V-A structure, the $SU(2)_L$ gauge construction, and the neutral-current and Higgs discoveries, to the modern electroweak Lagrangian, in order to motivate the neutrino electromagnetic form factors, and in particular the charge radius, as a natural and gauge-invariant extension of that framework rather than an ad hoc addition. Within that formalism we detailed the Dirac and Weyl cases for the NCR, and we clarified the numerical relation between the two conventions most commonly found in the literature, those of Bernab\'eu et al. and of Cadeddu et al., which differ by a factor of $-2$ and have historically been a source of confusion when comparing bounds across experiments.

We then derived, and made explicit, how a nonzero NCR modifies the effective proton vector coupling $g_p^V$ that enters the CE$\nu$NS cross section, inducing a calculable shift that is degenerate, at leading order, with a shift in $\sin^2\theta_W$ itself. This degeneracy is the central technical difficulty of the analysis, and it motivated the two complementary strategies pursued in this thesis: constraining $\sin^2\theta_W$ under the assumption of a SM NCR, and constraining $\langle r_{\nu_e}^2\rangle$ under the assumption of the SM weak mixing angle, for a family of realistic experimental configurations. The event rate was computed from the convolution of the Huber-Mueller reactor antineutrino flux model with the CE$\nu$NS differential cross section on $^{40}$Ar, for a detector placed 3 m from the core of the TRIGA Mark III reactor at ININ, near Mexico City, a compact and accessible neutrino source well suited to a staged, moderate-scale detector program.

A significant part of the analysis was devoted to the treatment of the detector response through the Lindhard quenching factor for argon, including the correct Jacobian treatment required to invert the relation between nuclear recoil energy and observed ionization energy. This step is not a minor technical detail: our results show that including the realistic quenching relation, rather than assuming an idealized unit response, degrades the projected precision on both $\sin^2\theta_W$ and $\langle r_{\nu_e}^2\rangle$ by a factor of two to three, since the low-energy recoils where the parameter dependence is most pronounced are precisely those most affected by quenching near threshold. We computed the projected $1\sigma$ sensitivities from a $\Delta\chi^2$ statistic, without correlations, for detector masses ranging from 10 to 200 kg and exposure times of 100 and 200 days. As expected from Poisson statistics, the sensitivity improves close to $1/\sqrt{N}$ with increasing exposure and close to an order of magnitude across the full mass range considered, with a 200 kg detector run for 200 days projected to reach a precision of order $1$-$2\%$ on $\sin^2\theta_W$ and of order $2\times10^{-33}\ \text{cm}^2$ on $\langle r_{\nu_e}^2\rangle$ even under the realistic quenching scenario. These figures place a reactor-based LAr array in a competitive, if not leading, position relative to existing low-energy determinations of the weak mixing angle, and they establish a concrete, testable target for the electron neutrino charge radius using a source and detector technology already available at a national laboratory scale.

Taken together, these results confirm that a modest, single-technology detector array is sufficient to place meaningful, simultaneous constraints on $\sin^2\theta_W$ and $\langle r_{\nu_e}^2\rangle$ from reactor CE$\nu$NS, provided that the quenching response is modeled with care and that exposure and target mass are pushed toward the upper end of what is technically feasible. The analysis presented here is, by design, a first approximation: it does not yet include full covariance between $\sin^2\theta_W$ and the NCR, nor does it explore additional flavor structure or physics beyond the Standard Model. Breaking the residual degeneracy between these two parameters, and extending the sensitivity study to the muon and tau neutrino charge radii and to non-standard neutrino interactions, is the natural continuation of this program, and it is precisely the direction outlined in the doctoral protocol that follows this thesis, which proposes a systematic, multi-flavor and multi-experiment study of the neutrino charge radius across SBC, CONUS+, IsoDAR, and DUNE. Each of these experiments probes a different combination of flux, target, and baseline, and a combined treatment across them, informed by the methodology developed here for the reactor-LAr case, offers a realistic path toward disentangling the weak mixing angle from the neutrino electromagnetic structure and toward constraining genuinely new physics in the neutrino sector.